% Compiler avec XeLaTeX ou LuaLaTeX
\documentclass[11pt,a4paper]{article}

\usepackage{amsmath,amssymb}
\usepackage{fontspec}
\usepackage{unicode-math}
\usepackage{polyglossia}
\setmainlanguage{french}
\usepackage{xcolor}
\usepackage{fancyhdr}
\usepackage[margin=2cm,headheight=16pt,headsep=22pt]{geometry}
\usepackage{enumitem}

% Police principale et police mathématique : Libertinus
\IfFontExistsTF{Libertinus Serif}
  {\setmainfont{Libertinus Serif}}
  {\IfFontExistsTF{Latin Modern Roman}{\setmainfont{Latin Modern Roman}}{\setmainfont{lmroman10-regular.otf}}}

\IfFontExistsTF{Libertinus Math}
  {\setmathfont{Libertinus Math}}
  {\IfFontExistsTF{Latin Modern Math}{\setmathfont{Latin Modern Math}}{\setmathfont{latinmodern-math.otf}}}

\linespread{1.07}
\setlength{\parindent}{0pt}
\AtBeginDocument{\setlength{\parindent}{0pt}}
\setlength{\parskip}{0.35em}

% En-tête
\pagestyle{fancy}
\fancyhf{}
\lhead{Professeur de Mathématiques}
\rhead{Seconde}
% Pagination : « Page X / Y » centrée en pied de page
\usepackage{lastpage}
\cfoot{Page \thepage\ / \pageref{LastPage}}
\renewcommand{\headrulewidth}{0.5pt}
\renewcommand{\headrule}{\vskip 5pt\hrule width\headwidth height\headrulewidth \vskip-\headrulewidth}

% Configuration des listes
\setlist[enumerate,1]{itemsep=0.45em, topsep=0.3em, label=\textbf{\arabic*.}}
\setlist[enumerate,2]{itemsep=0.25em, topsep=0.2em, label=\textbf{\alph*.}}

% Commande exercice
\newcommand{\exo}[2]{%
  \vspace{1.5em}%
  \noindent{\large\textbf{Exercice #1. -- #2}}%
  \par\vspace{0.35em}%
  \hrule height 0.5pt%
  \vspace{0.7em}%
}

\newcommand{\partie}[1]{%
  \vspace{0.9em}%
  \noindent\textbf{#1}%
  \par\vspace{0.3em}%
}

\begin{document}

\begin{center}
  {\LARGE \textbf{Exercice d'entraînement : Racines carrées et Identités remarquables}}\\[0.5em]
  {\large Classe de Seconde}
\end{center}

% ---------------------- EXERCICE 1 ----------------------
\exo{1}{Calculs et simplifications algébriques}

\partie{Partie A -- Simplification de radicaux}
Écrire chacun des nombres suivants sous la forme $a\sqrt{b}$ où $a$ et $b$ sont des entiers, $b$ étant le plus petit possible.
\begin{enumerate}
  \item $A = \sqrt{75} - 2\sqrt{12} + 3\sqrt{27}$
  \item $B = \sqrt{50} \times \sqrt{8}$
\end{enumerate}

\partie{Partie B -- Utilisation des identités remarquables}
Développer et simplifier les expressions suivantes en détaillant les étapes :
\begin{enumerate}
  \item $C = (3 + \sqrt{5})^2$
  \item $D = (2\sqrt{3} - 4)(2\sqrt{3} + 4)$
  \item $E = (\sqrt{7} - 2)^2 + (\sqrt{7} + 2)(\sqrt{7} - 2)$
\end{enumerate}

% ---------------------- EXERCICE 2 ----------------------
\exo{2}{Étude d'une expression littérale et application géométrique}

On considère l'expression numérique $F = (2\sqrt{2} - 1)^2 + 4\sqrt{2}$.

\begin{enumerate}
  \item Développer et réduire l'expression $F$ en utilisant l'identité remarquable adaptée. Écrire le résultat sous la forme $a + b\sqrt{2}$ où $a$ et $b$ sont des entiers.
  \item Soit un rectangle de dimensions $L = 2\sqrt{2} - 1$ et $l = \sqrt{2}$. 
  \begin{enumerate}
    \item Exprimer son aire exacte sous forme développée et simplifiée.
    \item Calculer la longueur exacte de sa diagonale $d$, puis donner la valeur simplifiée de $d^2$.
  \end{enumerate}
\end{enumerate}

\vspace{2em}
\hrule
\vspace{1em}
\noindent{\large\textbf{Corrigé détaillé}}

\partie{Exercice 1}
\textbf{Partie A :}
\begin{enumerate}
  \item $A = \sqrt{25 \times 3} - 2\sqrt{4 \times 3} + 3\sqrt{9 \times 3} = 5\sqrt{3} - 2 \times 2\sqrt{3} + 3 \times 3\sqrt{3} = 5\sqrt{3} - 4\sqrt{3} + 9\sqrt{3} = 10\sqrt{3}$.
  \item $B = \sqrt{50 \times 8} = \sqrt{400} = 20$.
\end{enumerate}

\textbf{Partie B :}
\begin{enumerate}
  \item $C = 3^2 + 2 \times 3 \times \sqrt{5} + (\sqrt{5})^2 = 9 + 6\sqrt{5} + 5 = 14 + 6\sqrt{5}$.
  \item $D = (2\sqrt{3})^2 - 4^2 = 4 \times 3 - 16 = 12 - 16 = -4$.
  \item $E = \left[(\sqrt{7})^2 - 2\times\sqrt{7}\times 2 + 2^2\right] + \left[(\sqrt{7})^2 - 2^2\right] = (7 - 4\sqrt{7} + 4) + (7 - 4) = (11 - 4\sqrt{7}) + 3 = 14 - 4\sqrt{7}$.
\end{enumerate}

\partie{Exercice 2}
\begin{enumerate}
  \item $F = (2\sqrt{2})^2 - 2 \times 2\sqrt{2} \times 1 + 1^2 + 4\sqrt{2} = 8 - 4\sqrt{2} + 1 + 4\sqrt{2} = 9$.
  \item 
  \begin{enumerate}
    \item $\text{Aire} = L \times l = (2\sqrt{2} - 1)\sqrt{2} = 2(\sqrt{2})^2 - \sqrt{2} = 2(2) - \sqrt{2} = 4 - \sqrt{2}$.
    \item D'après le théorème de Pythagore, $d^2 = L^2 + l^2 = (2\sqrt{2} - 1)^2 + (\sqrt{2})^2 = (9) + 2 = 11$. Donc $d = \sqrt{11}$.
  \end{enumerate}
\end{enumerate}

\end{document}
